% !TEX root = AImath.v4.tex
\chapter{ 引言 什么是智能？}
 \begin{introduction}
  \item Definition of Theorem
  \item Ask for help
  \item Optimization Problem
  \item Property of Cauchy Series
  \item Angle of Corner
\end{introduction}


这一章将带领你从哲学、科学和工程的角度出发，追问“什么是智能”。接下来的章节中，我们将深入探讨数学在人工智能中的具体应用，展示数学如何帮助我们创造出越来越强大的智能系统。

\section{什么是智能与人工智能}

{\textbf{智能(intelligence)，如今已经成为现代生活中很常见的一个词，比如智能手机、智能家居、智能驾驶等。在不同使用场合中，智能的含义也不太一样。比如“智能手机”中的“智能”一般是指由计算机控制并具有某种智能行为的意思。这里的“计算机控制”+“智能行为”隐含了对人工智能的简单定义。}}

%{简单地讲，人工智能(Artificial Intelligence，AI)就是让机器具有人类的智能，这也是人们长期追求的目标。这里关于什么是“智能”并没有一个很明 确的定义，但一般认为智能(或特指人类智能)是知识和智力的总和，都和大脑的思维活动有关。人类大脑是经过了上亿年的进化才形成了如此复杂的结构，但至今仍然没有完全了解。虽然随着神经科学、认知心理学等学科的发展，人们对大脑的结构有了一定程度的了解，但对大脑的智能究竟怎么产生的还知道的很少。我们并不理解大脑的运作原理，以及如何产生意识、情感、记忆等功能。因此，通过“复制”一个人脑来实现人工智能在目前阶段是不切实际的。}

%智能，或许是我们人类最自豪的天赋。
智能并非人类独有，许多动物，如海豚、鸟类、猩猩，甚至章鱼，都表现出高度的学习能力、问题解决能力以及某种形式的社交智能，但人类的智能，但人类的智能以其复杂性和多样性独树一帜，在抽象和创造性上尤为突出。人类具有抽象思维，创造了诸多符号系统，并基于此进行高度复杂的推理和想象。从数学公式到哲学思辨，人类通过语言、文字和符号，建立起了庞大的知识体系和文化遗产。我们可以预见未来、反思过去，并通过科学、技术、工具改变现实世界。正是这种智能上的深度和广度，推动了人类文明的发展。

智能是生命进化过程中形成的一种适应性特征。但随着人工智能的发展，我们正在见证一种新型智能的崛起。在20世纪中期，计算机的发明标志着一种全新“智能”形式的诞生。计算机，这种基于二进制系统和算法逻辑的机器，能够完成许多过去被认为只有人类才能做到的任务，如图像识别、语言翻译、甚至写作创作。机器可以通过算法模拟复杂的推理过程，甚至在某些领域超越人类的能力——它们能够比人类更快地进行大规模计算、更高效地分析海量数据。

这促使我们重新思考：究竟什么是“智能”？它是基于生物特性的一种特有思维能力，还是某种可以通过数学和逻辑建模来实现的计算逻辑？

对于智能的定义，在不同学科中有着不同的阐释。
在心理学中，智能被视为一种个体解决问题、学习和适应环境的能力。
在生物学中，智能更多与生物体的行为、神经系统的复杂性相关。
而在计算机科学和人工智能领域，智能则被看作是一种能够在复杂环境中执行任务并达成特定目标的能力。
从人工智能的角度看，智能可以被定义为“机器执行通常需要人类智能才能完成的任务的能力”。这些任务包括视觉识别、自然语言处理、学习、推理、决策等。然而，计算机和人类智能之间仍然存在着巨大的差异，特别是在理解、情感和创造力等方面。人工智能的核心在于如何用数学工具来模仿、简化甚至扩展这些认知能力。

机器“智能”是否真的与人类的智能一样？人类的智能伴随着进化而来，经过了数百万年的演变和适应，它不仅仅是一种计算能力，更是情感、直觉、创造力和推理能力的综合体。相对来说，人工智能的发展历史尚短，但其进化速度却让人惊叹。

机器智能的本质在于其基于算法和数据的计算逻辑，而人类智能则包含了情感、意识、创造力等更为复杂和难以定义的元素。智能不仅仅是任务的完成能力，更涉及到意识、理解和体验。也许机器可以在规则内完成任务，优化目标，但它们仍然无法像人类那样拥有感知、情感或独立的创造性思维。尽管如此，在这场由人类主导的智能革命中，我们依然是规则的制定者和方向的掌控者。这也是为什么，理解数学思维如何推动人工智能，能够帮助我们更好地把握未来智能发展的关键脉络。

\subsection*{数学在智能中的角色}

自古以来，数学不仅是解决问题的工具，更是智能表达与实现的基础语言。数学的力量在于它能够抽象出复杂问题的本质，提供精确的描述和推理方法。

人工智能中的诸多任务，如模式识别、数据分类、预测等，背后都依赖数学模型进行量化和求解。例如，图像识别问题看似复杂，但可以被归纳为一个线性代数问题：图像被分解为像素点，每个像素点通过向量和矩阵的运算来进行处理。同样，神经网络的训练也依赖于微积分中的优化技术，以最小化损失函数，找到最优参数组合。

\subsection*{机器智能的进化：从计算到学习}

数学在推动人工智能从简单计算迈向学习与适应中起到了关键作用。最初，计算机的智能主要体现在通过预先编写的规则来执行任务，模拟人类的某些智能行为。随着大数据和深度学习技术的出现，人工智能开始从数据中“学习”，即所谓“机器学习”。这是一场范式的转变：机器智能从依赖人工预定义的规则，转变为通过模型的不断训练和优化进行自我学习。

在这场范式的转变中，数学再次发挥了至关重要的作用。统计学、线性代数、概率论等数学工具，使机器能够从复杂的海量数据中提取有价值的模式和信息，并进行合理决策。

机器学习依赖于统计学、线性代数、概率论等数学理论。通过这些工具，机器能够从海量数据中提取模式并作出决策，具备了学习和自我优化的能力。

这种数学思维下的机器学习能力，使得智能不仅仅是计算结果的简单输出，而是一种能够学习、适应和自我优化的动态过程。这场智能的进化，无疑奠定了人工智能发展的新高度。

\subsection*{数学思维：解码智能的钥匙}

人工智能是模拟、延伸和扩展人类智能的理论、方法、技术及其应用系统的综合学科。AI作为一门跨学科领域，从多个学科中汲取灵感与启发，融合了不同领域的知识与技术。

数学思维不仅仅是对现有工具的应用，更代表着一种解码世界的方式。数学能够帮助我们从表象中剥离出系统的本质，理清复杂事物之间的逻辑关系。在人工智能领域，数学思维是帮助我们更好理解并发展智能的核心钥匙。
人工智能的本质是一种通过数学和计算实现的智能，它利用算法、数据处理和计算资源模拟生物智能的行为。
为了理解智能如何通过数学实现，我们可以将其分解为几个关键组成部分：
\begin{itemize}
\item {\bf 感知与反应：} 任何智能体首先需要能够感知环境，并对环境做出适当的反应。这是智能的基础，也是进化中所有生物生存的关键。数学在感知数据的处理上扮演重要角色，例如通过卷积神经网络处理视觉输入。

\item  {\bf 学习与记忆：} 智能体不仅是对环境做出反应，还包括从经验（数据）中学习，，储存经验用于未来决策。

\item  {\bf 推理与决策：} 智能体需要能够进行推理，从已有信息中得出新的结论，从而做出合理的决策。这离不开统计推断、优化等数学方法。

\item  {\bf 创造与适应：} 更高级的智能不仅仅是反应式的，而是能够创造性地应对未知的挑战，并在复杂的环境中灵活适应。
\end{itemize}
人工智能（AI）在这些维度上逐步模仿人类的智能，但它的本质是通过数学和计算实现的。AI的智能是通过算法、数据处理以及计算资源，以我们在生物智能中所见的行为为目标进行的复杂优化。

、\section{鹦鹉与乌鸦：何种智能?}

为了更好地理解智能，我们可以借用一个生动的比喻——鹦鹉与乌鸦。

\subsection*{鹦鹉式智能：模仿而不理解}
鹦鹉是模仿高手，能够复制人类的语言，甚至在特定情境下做出看似有意识的反应。然而，鹦鹉并不真正理解它模仿的内容，它只是重复我们的话语，而无法将这些词汇与现实世界中的人、物或情景建立联系。这种行为类似于一些现有的人工智能系统，特别是那些专注于模仿和执行特定任务的AI模型。这类AI可以通过大量数据训练，在特定任务中表现优异，但它们并不真正理解背后的语义或意图。

这种“鹦鹉式智能”可以类比为{\bf 弱AI}或{\bf 窄AI}。它们擅长于特定任务，但缺乏广泛的推理和创造能力。当前的许多AI模型，如语言生成或图像识别，都是在这一框架下运作的——它们通过模式识别和数据驱动的方式执行任务，却没有真正的自主意识或深度理解。都是在这一框架下运作的——它们通过模式识别和数据驱动的方式执行任务，却没有真正的自主意识或深度理解。

\subsection*{乌鸦式智能：推理与创造}

乌鸦向来被认为是最聪明的鸟类之一。与鹦鹉不同，乌鸦不仅擅长模仿，还展现出复杂推理能力、问题解决能力以及创造性的智能。乌鸦是少数能够制造工具的鸟类，能够通过创新思维应对物理世界中的挑战。
我们从小都听过各种“乌鸦喝水”的故事，而国外科学家也曾做过一系列“乌鸦喝水”的变化实验。在这些实验中，研究人员使用了各种不同形状的瓶子和不同密度的物体，观察乌鸦如何应对。令人惊讶的是，乌鸦总能在几次尝试后，快速识别并选择那些能够最快让水位上升的物体。这些实验不仅证实了乌鸦对物理规律的深刻理解，还展示了它们高度的认知能力和推理能力，以及面对复杂问题时的创造性和灵活性。
另一个著名的例子是，乌鸦会将坚果放在马路上，利用来往车辆碾压坚果壳，并观察红绿灯、斑马线等信号来推理交通规则，最终推理出红绿灯与车辆停靠及行人停走之间的复杂因果关系，从而安全地过马路，吃到被碾碎的坚果肉。这是一种远远超出简单反应的行为，展示了乌鸦对物理世界和人类行为之间复杂因果链的深刻理解。

乌鸦的这一行为没有经过大量数据训练或监督学习，它通过少量的尝试与观察，直接推理出解决问题的方法。乌鸦智能正是我们期望的真正的智能---完全自主的智能：观察、感知、认知、学习、推理、执行。这种智能可以比喻为{\bf 强AI}或{\bf 通用AI}，具备自主学习、推理和适应能力，能够在不同的环境中灵活应对各种任务。

功耗极低：乌鸦头不到人脑的1%大小。人脑功耗大约是10-25瓦，乌鸦头只有0.1-0.2瓦，不需要核动力发电、冷却。

乌鸦的智能不仅令人惊叹，功耗也极低。人类大脑的功耗大约在10-25瓦之间，乌鸦的大脑尺寸约为人类的1\%，功耗约为0.1-0.2瓦。这意味着乌鸦无需复杂的能源供应系统或冷却设备，就能完成复杂的推理和问题解决任务。相比之下，现代AI系统需要高性能计算设备和大量能源，显示了生物智能的高效性。

\subsection*{智能的层次：从鹦鹉到乌鸦}
“鹦鹉”与“乌鸦”的比喻在人工智能领域有着深刻的象征意义，揭示了{\bf 智能的不同层次}：
\begin{itemize}
\item {\bf 鹦鹉式智能：} 擅长特定任务，通过数据驱动的模式识别和模仿来实现，针对某些垂直应用有效，但缺乏深度理解和推理能力。
\item {\bf 乌鸦式智能：} 具备复杂的推理和创造性，能够自主学习、适应环境，展现出高水平的智能灵活性。
\end{itemize}

这个比喻通过比较鹦鹉的模仿性智能和乌鸦的推理性智能，说明了AI智能的不同层次。

\subsection*{AI的未来：向“乌鸦式智能”迈进}
现有的AI系统，即使是最先进的言模型和图像识别系统，都还处于“鹦鹉式智能”阶段。这些系统依赖于预定义规则、大量数据、大量计算资源来完成特定任务，但缺乏真正的推理和创造能力。虽然这类AI能够生成看似“智能”的对话，但它们的智能基于模式识别和数据驱动，与鹦鹉能模仿人类发音但不理解其语义一样。即使是像ChatGPT这样的先进语言模型，在模仿人类语言并生成非常连贯对话上表现几可乱真，但它并不具备关于“意义”的主观理解，无法感知、体验或有意识地推理，只是基于统计模式和相关性根据输入生成最有可能符合上下文的回应。因此，大语言模型更多是“聪明的猜测”，而非真正的理解。这与人类的理解方式存在本质差异。


AI的未来发展目标是朝着“乌鸦式的智能”迈进，不仅能够模仿，还能推理、学习、适应，在跨任务和多变的环境中展现出灵活的创造力。实现与生物智能相媲美的灵活性和创造力，正是“强AI”或“通用AI”所要追求的方向。


人工智能:		模拟、延伸和扩展人的智能   的理论、方法、技术及应用系统


\nocite{*}
%\printbibliography[segment=\therefsegment, heading=bibintoc, title=\ebibname]

\newpage
